\documentclass[10pt]{article}

\usepackage[margin=23mm]{geometry}
\usepackage{amsmath,amssymb}
\usepackage{booktabs}
\usepackage{array}
\usepackage{enumitem}
\usepackage{xcolor}
\usepackage{hyperref}
\usepackage{fancyhdr}

\definecolor{linkblue}{HTML}{1F5A7A}
\definecolor{softblue}{HTML}{EAF3F7}
\definecolor{softgrey}{HTML}{F3F3F3}
\hypersetup{
  colorlinks=true,
  linkcolor=linkblue,
  urlcolor=linkblue,
  citecolor=linkblue,
  pdftitle={Parametric Brownian Bootstrap for the Full C. elegans Lineage},
  pdfauthor={Embryogenesis Pareto analysis project}
}

\setlist{nosep,leftmargin=5mm}
\setlength{\parindent}{0pt}
\setlength{\parskip}{5pt}
\setlength{\tabcolsep}{5pt}
\renewcommand{\arraystretch}{1.15}

\pagestyle{fancy}
\fancyhf{}
\lhead{Full-tree Pareto analysis}
\rhead{Brownian parametric bootstrap}
\cfoot{\thepage}
\setlength{\headheight}{14pt}

\newcommand{\R}{\mathbb{R}}
\newcommand{\E}{\mathbb{E}}
\newcommand{\Var}{\operatorname{Var}}
\newcommand{\Cov}{\operatorname{Cov}}
\newcommand{\Normal}{\mathcal{N}}
\newcommand{\Tobs}{T_{\mathrm{obs}}}
\newcolumntype{L}[1]{>{\raggedright\arraybackslash}p{#1}}

\title{\textbf{A Parametric Brownian Bootstrap for the\\
Full \textit{C. elegans} Lineage}\\[2mm]
\large Technical rationale, model, algorithm, validation, and interpretation}
\author{Embryogenesis Pareto analysis project}
\date{1 September 2026; superseded for Figure 8 on 9 September 2026}

\begin{document}
\maketitle

\section*{Historical specification: superseded for Figure 8}
The shared-tracking-time model documented below is retained for historical
comparison. Figure 8 now uses the model specified in
\texttt{separate\_clock\_reference.tex}: tracking-time spatial covariance
and protein covariance per canonical lineage transition, with independent
blocks. Aggregate protein states are not aligned to cutoff-truncated tracking
endpoints, so division by short tracking durations can inflate their fitted
rate. Adoption on 9 September 2026 was based on this measurement-clock
mismatch, not on tuning the reference to observed costs. Statements below
about the proposed Figure 8 reference refer to the superseded specification.

\begin{abstract}
This historical note specifies the former Brownian-motion parametric bootstrap for evaluating the
travel and protein-state edge costs of the measured \textit{C. elegans}
lineage.  It is written for the current full-tree publication configuration:
four measured roots, 500 internal cells, 504 terminal cells, 1,000 evaluated
edges, three spatial coordinates, and 20 z-scored protein features.  The key
distinction is between \emph{ancestral-state reconstruction}, which infers
unobserved internal states conditional on terminal observations, and a
\emph{generative null}, which simulates complete observable lineages and
compares exactly the same statistics with the natural lineage.  The latter is
the appropriate role for a Brownian reference distribution in Figure~8.
\end{abstract}

\setcounter{tocdepth}{1}
\tableofcontents
\newpage

\section{The scientific question}

The full-tree analysis observes a state vector for every measured cell,
\[
  \mathbf{x}_v =
  \bigl(x_v,y_v,z_v,p_{v1},\ldots,p_{v20}\bigr)^\top \in \R^{23},
\]
where the first three entries are spatial coordinates and the remaining 20 are
z-scored protein features.  The natural forest contains four earliest measured
roots and a common set $E$ of 1,000 scored parent--child edges.

The proposed null asks:

\begin{quote}
If cell states changed by neutral, constant-rate Gaussian diffusion along the
same measured branches, what total travel and protein-state costs would be
observed?
\end{quote}

This is a model-adequacy question.  It holds the measured forest and branch
durations fixed, generates complete alternative state histories under a fitted
Brownian model, and evaluates the same two statistics used for the natural
lineage.  Simulation-based tests on a fixed phylogeny have a long history
\cite{garland1993}; modern phylogenetic Monte Carlo and model-adequacy
frameworks emphasize fitting an explicit evolutionary model and checking its
predictions on the particular tree and statistic of interest
\cite{boettiger2012,pennell2015}.

\subsection{What this null does and does not test}

\begin{itemize}
  \item It tests neutral \emph{state dynamics} on the fixed natural forest.
  \item It does not randomize lineage topology, parent identities, branch
        durations, measurement availability, or the number of scored edges.
  \item It therefore complements, rather than replaces, the assignment,
        cousin-shuffle, and reconstruction nulls in the Pareto analysis.
  \item A result outside the bootstrap distribution rejects this particular
        Brownian description; it does not by itself identify selection or a
        unique biological mechanism.
\end{itemize}

\section{Why the former method answers a different question}

The historical implementation begins with terminal states, estimates
Brownian rates from them, and infers internal states conditional on those
terminals.  In probability notation it targets
\[
 p\!\left(
   \mathbf{X}_{\mathrm{internal}}
   \mid \mathbf{X}_{\mathrm{terminal}},\mathcal{T},\widehat{\boldsymbol\Sigma}
 \right).
\]
This is ancestral-state reconstruction.  It is appropriate when the internal
states are missing and the goal is to predict them.  Studies with independently
observed ancestral nodes use those observations to evaluate the reconstructions
rather than treating reconstructed values as new neutral data
\cite{polly2001}.

In the present experiment, however, the 500 internal cells are observed.  If
they are discarded and replaced by terminal-conditioned estimates, those
estimates act as smooth interpolants.  Comparing their edge lengths with the
irregular measured trajectory gives the reconstruction an intrinsic smoothness
advantage.  It answers ``what smooth ancestral history explains these tips?''
rather than ``what complete lineage would neutral Brownian dynamics generate?''

\subsection{Three additional problems in the historical code}

\begin{enumerate}
  \item \textbf{Identity indexing.} Tree IDs encode topology; lineage IDs index
        measurement arrays.  None of the 504 terminal tree IDs in the current
        configuration equals its biological measurement-row ID.  Directly
        indexing the data by tree ID assigns the wrong cell states to every tip.

  \item \textbf{Independent node draws.} Marginal ancestral distributions were
        sampled independently.  A Brownian parent and child are positively
        correlated.  For adjacent states $P$ and $C$,
        \[
          \Var(C-P)=\Var(C)+\Var(P)-2\Cov(C,P).
        \]
        Dropping the positive covariance inflates edge variation.

  \item \textbf{Unequal scopes.} The former sampler summed 1,006 represented
        edges, including early construction edges.  Natural and heuristic
        costs use the 1,000 edges below the four measured roots.
\end{enumerate}

The old 10-PC experiment illustrates the combined effect.  Its natural
expression cost is 93,661.  The historical calculation with incorrect
identity indexing yields about 106,857, reproducing the old plotted cloud.
Correct identity mapping lowers the same marginal-node method to about 70,338;
exact joint conditioning lowers it to about 67,464.  The low value is not a
numerical accident: it is the expected smoothness of terminal-conditioned
ancestral reconstruction.

\section{Generative Brownian model}

Let $e=(u\rightarrow v)$ denote a directed edge from parent $u$ to child $v$,
and let $t_e>0$ be its tracking-time duration.  A multivariate Brownian model
assumes
\begin{equation}
  \mathbf{x}_v = \mathbf{x}_u + \boldsymbol\epsilon_e,
  \qquad
  \boldsymbol\epsilon_e \sim
  \Normal\!\left(\mathbf{0},t_e\boldsymbol\Sigma\right),
  \label{eq:brownian}
\end{equation}
with independent increments on distinct edges.  The $23\times23$ positive
semidefinite matrix $\boldsymbol\Sigma$ describes change per unit branch time,
including feature variances and cross-feature covariances.  Under Brownian
motion, variance accumulates in direct proportion to elapsed branch time, the
central assumption behind independent contrasts and phylogenetic generalized
least squares \cite{felsenstein1985,freckleton2012}.

Equation~\eqref{eq:brownian} is a zero-drift model: the expected increment is
zero.  A directional extension is
\[
  \boldsymbol\epsilon_e \sim
  \Normal(t_e\boldsymbol\mu,t_e\boldsymbol\Sigma),
\]
but drift should be added only as a prespecified robustness analysis.  With
many correlated protein features, an unconstrained drift vector can absorb
developmental trends and materially change the biological null.

\subsection{Role of the topology}

Simulating absolute node states requires the tree: a terminal state is the sum
of all increments along its root-to-tip path, and related cells share early
increments.  This produces the standard phylogenetic covariance structure.

For the particular statistic used here---a sum of norms of individual edge
increments---the distribution depends on the collection of branch durations
and the fitted increment covariance, but not on how edges are connected.  The
topology becomes essential again for terminal covariances, clade summaries,
subtree statistics, or any statistic involving nonadjacent cells.  Stating
this explicitly prevents overinterpreting the bootstrap as a topology test.

\section{Parameter estimation from the measured forest}

For each of the same 1,000 scored edges, calculate the observed increment
\[
  \boldsymbol\Delta_e = \mathbf{x}_v-\mathbf{x}_u
\]
and its branch-standardized form
\[
  \mathbf{z}_e = \frac{\boldsymbol\Delta_e}{\sqrt{t_e}}.
\]
Under Equation~\eqref{eq:brownian}, the $\mathbf{z}_e$ are independent draws
from $\Normal(\mathbf{0},\boldsymbol\Sigma)$.  The zero-drift maximum-likelihood
estimate is therefore
\begin{equation}
  \widehat{\boldsymbol\Sigma}
    = \frac{1}{|E|}\sum_{e\in E}
      \frac{\boldsymbol\Delta_e\boldsymbol\Delta_e^\top}{t_e}.
  \label{eq:sigmahat}
\end{equation}

This uses the internal observations rather than treating them as missing.
Because $|E|=1{,}000$ and the state dimension is 23, the full empirical
covariance is estimable, but its eigenvalues and condition number must still be
reported.

\subsection{Covariance specifications to compare}

\begin{center}
\begin{tabular}{L{28mm}L{50mm}L{72mm}}
\toprule
Specification & Structure & Interpretation \\
\midrule
Diagonal & 23 variances & Reproduces the historical feature-independence
assumption; simplest but discards correlated protein changes. \\
Two-block & Full $3\times3$ spatial and $20\times20$ protein blocks & Preserves
correlation within each objective while preventing cross-objective covariance
from influencing the generative model. \\
Full & Full $23\times23$ covariance & Preserves spatial--protein coupling as
well as within-objective covariance; preferred primary model if numerically
stable and scientifically intended. \\
Shrinkage & Regularized full covariance & Useful if eigenvalues are unstable;
the shrinkage rule and intensity must be fixed and reported. \\
\bottomrule
\end{tabular}
\end{center}

The diagonal, two-block, and full specifications form an important robustness
analysis.  They should not be selected after inspecting which one makes the
natural lineage look most extreme.

\section{Bootstrap algorithm}

Let $B$ be the number of replicates, preferably at least 10,000 for a
publication null cloud and stable tail estimates.  Compute a matrix square root
$\mathbf{L}$ such that
\[
  \mathbf{L}\mathbf{L}^\top=\widehat{\boldsymbol\Sigma}.
\]
Cholesky decomposition is appropriate for a positive-definite estimate;
eigendecomposition with documented clipping or shrinkage is safer near the
positive-semidefinite boundary.

For replicate $b=1,\ldots,B$:

\begin{enumerate}
  \item Fix the four root states at their measured values.  Root values do not
        affect edge-cost totals, but retaining them permits complete simulated
        node-state output and later topology-aware diagnostics.
  \item Traverse from roots toward tips.  For every scored edge $e$, draw
        $\mathbf{g}_e^{(b)}\sim\Normal(\mathbf{0},\mathbf{I}_{23})$ and set
        \[
          \boldsymbol\Delta_e^{(b)}
          = \sqrt{t_e}\,\mathbf{L}\mathbf{g}_e^{(b)},
          \qquad
          \mathbf{x}_v^{(b)}
          = \mathbf{x}_u^{(b)}+\boldsymbol\Delta_e^{(b)}.
        \]
  \item Evaluate exactly the natural-lineage statistics:
        \begin{align}
          C_{xyz}^{(b)} &=
            \sum_{e\in E}\left\|\boldsymbol\Delta_{e,1:3}^{(b)}\right\|_2,
            \\
          C_{protein}^{(b)} & =
            \sum_{e\in E}\left\|\boldsymbol\Delta_{e,4:23}^{(b)}\right\|_2.
        \end{align}
  \item Store costs, the seed or replicate index, and any model-adequacy
        summaries.  Do not store a million full node-state tensors unless a
        later statistic requires them.
\end{enumerate}

\subsection{Reference pseudocode}

\begin{verbatim}
# Fit on exactly the evaluated 1,000 edges
delta[e] = state[child[e]] - state[parent[e]]
z[e]     = delta[e] / sqrt(branch_time[e])
Sigma    = mean(outer(z[e], z[e]) for e in evaluated_edges)
L        = matrix_square_root(Sigma)

for b in range(B):
    travel = 0
    protein = 0
    simulated_state[root_ids] = observed_state[root_ids]

    for e in root_to_tip_edge_order:
        increment = sqrt(branch_time[e]) * L @ normal_vector(23)
        simulated_state[child[e]] = simulated_state[parent[e]] + increment
        travel  += norm(increment[:3])
        protein += norm(increment[3:])

    save(b, travel, protein)
\end{verbatim}

For edge-additive costs alone, the absolute-state assignments may be omitted
and independent increments accumulated directly.  Retaining the traversal is
recommended in the reference implementation because it enables tests of
terminal covariance and subtree behavior.

\section{Inference from the bootstrap distribution}

The observed statistics are
\[
  C_{xyz}^{obs}=4{,}465.55,
  \qquad
  C_{protein}^{obs}=3{,}102.41.
\]
If low cost is the prespecified direction of interest, finite-sample empirical
$p$-values are
\begin{align}
 p_{xyz}^{low}
  &= \frac{1+\sum_{b=1}^{B}
       \mathbf{1}\!\left[C_{xyz}^{(b)}\le C_{xyz}^{obs}\right]}{B+1},\\
 p_{protein}^{low}
  &= \frac{1+\sum_{b=1}^{B}
       \mathbf{1}\!\left[C_{protein}^{(b)}\le C_{protein}^{obs}\right]}{B+1}.
\end{align}
The $+1$ correction prevents an estimated probability of exactly zero.  Report
the bootstrap mean, standard deviation, central interval, empirical percentile,
Monte Carlo sample size, seed, and covariance specification alongside each
$p$-value.

Because the two costs are evaluated jointly, Figure~8 should also display the
two-dimensional cloud.  A joint extremeness measure must be chosen before
examining results.  Reasonable choices include a lower-left orthant probability
or a Mahalanobis discrepancy in the bootstrap cost plane.  Neither should be
silently substituted for the scientifically clearer pair of marginal results.

\subsection{Parameter uncertainty}

The primary plug-in bootstrap simulates from
$\widehat{\boldsymbol\Sigma}$.  Parameter uncertainty can be assessed by an
outer bootstrap over standardized measured increments, by a Bayesian posterior
over $\boldsymbol\Sigma$, or by a nested parametric bootstrap.  These are
valuable robustness checks, especially for tail probabilities, but they do not
repair a structurally inadequate Brownian model.  The model must first pass
basic adequacy diagnostics.

\section{Model-adequacy checks required before interpretation}

A small $p$-value means the observed statistic is unusual under the fitted
model.  It is interpretable only after asking why the model failed.  Pennell
et~al. emphasize that absolute adequacy is distinct from choosing the best
model among several inadequate alternatives \cite{pennell2015}.

The release analysis should include at least the following checks:

\begin{enumerate}
  \item \textbf{Standardized increments.} Examine
        $\mathbf{z}_e=\boldsymbol\Delta_e/\sqrt{t_e}$ for mean trends,
        skewness, heavy tails, and outliers.
  \item \textbf{Rate homogeneity.} Compare squared standardized increments by
        contraction round, developmental time, major lineage, and terminal
        versus internal edge.  Constant-rate Brownian motion predicts no
        systematic grouping.
  \item \textbf{Branch-time calibration.} Test whether squared increment
        magnitude is proportional to branch duration.  Tracking duration is a
        modeling choice, not automatically an evolutionary clock.
  \item \textbf{Feature covariance.} Report covariance eigenvalues, condition
        number, effective rank, and results under diagonal, two-block, and full
        models.
  \item \textbf{Posterior-predictive summaries.} In addition to total cost,
        compare maximum edge change, coefficient of variation of edge norms,
        layerwise costs, subtree costs, and terminal covariance with their
        simulated distributions.
  \item \textbf{Monte Carlo convergence.} Repeat with at least two seeds and
        show that means, intervals, and tail probabilities are stable.
\end{enumerate}

If the model fails because rates differ greatly across developmental layers,
a prespecified multi-rate Brownian model may be appropriate.  If state changes
show attraction to a developmental trajectory or optimum, an
Ornstein--Uhlenbeck-like model may be considered.  If increments are
heavy-tailed, a jump or robust diffusion model may be more credible.  These
models answer different biological questions and should be motivated by
diagnostics, not selected solely to move the null cloud.

\section{Preliminary numerical audit and expected figure behavior}

A 4,000-draw diagnostic run using the full $23\times23$ covariance estimated
from all 1,000 observed increments produced:

\begin{center}
\begin{tabular}{lrr}
\toprule
 & Travel cost & Protein-state cost \\
\midrule
Natural lineage & 4,465.55 & 3,102.41 \\
Brownian bootstrap mean & 4,831.30 & 5,241.96 \\
Brownian bootstrap SD & 74.16 & 38.00 \\
Mean / natural & 1.082 & 1.690 \\
\bottomrule
\end{tabular}
\end{center}

These values are an audit preview, not a release result.  They show why the new
null is conceptually more plausible than terminal-conditioned reconstruction:
the simulated lineages are complete observable histories, and the natural
lineage is compared with the same 1,000-edge statistic.  They also suggest
strong Brownian model inadequacy for protein-state cost.  Before interpreting
that separation as biological organization, the constant-rate and Gaussian
assumptions must be checked using the diagnostics above.

In Figure~8 coordinates standardized by the first-cousin null, this bootstrap
may lie far from the main Pareto region on the protein axis.  An inset is then
appropriate, but its distance should be described as a model-adequacy result,
not as evidence that Brownian motion is the uniquely correct neutral process.

\newpage
\section{Reproducibility and validation checklist}

The implementation is ready for publication use only if all of the following
are recorded and tested:

\begin{itemize}
  \item exact biological identity mapping for every parent and child;
  \item exactly four roots, 1,000 unique evaluated edges, and no early-anchor
        construction edges;
  \item positive, explicitly defined branch durations for every edge;
  \item the 3 spatial and 20 protein features in documented order and scale;
  \item covariance specification, estimator, regularization, eigenvalues, and
        condition number;
  \item random-number generator, seed, replicate count, and cache schema;
  \item deterministic agreement between vectorized and reference traversal
        implementations on a small seed;
  \item agreement of natural costs recomputed from the bootstrap edge manifest
        with 4,465.55 and 3,102.41;
  \item Monte Carlo convergence across seeds and replicate counts;
  \item model-adequacy summaries, not only the final two cost coordinates; and
  \item clear labeling as a parametric Brownian reference on a fixed topology,
        not an ancestral reconstruction or a topology-randomization null.
\end{itemize}

\section{Recommended interpretation}

The parametric bootstrap has a narrow, useful interpretation:

\begin{quote}
After calibrating a constant-rate multivariate Brownian process to the observed
full-tree increments, does that process reproduce the natural lineage's total
travel and protein-state edge costs on the same branch durations?
\end{quote}

Its advantages are scope matching, correct joint generation, explicit
assumptions, direct simulation of the plotted statistics, and the ability to
diagnose model failure.  Its limitations are equally important: the Brownian
assumption may be poor for a deterministic developmental lineage; tracking
time may not be the correct diffusion clock; feature scaling affects the
covariance; and a fixed-topology test cannot establish that the natural parent
assignment itself is special.

Used with those qualifications, the bootstrap is a defensible reference model.
Used without adequacy checks, it risks replacing one misleading phylogenetic
cloud with another.

\begin{thebibliography}{9}

\bibitem{felsenstein1985}
Felsenstein, J. (1985).
Phylogenies and the comparative method.
\textit{The American Naturalist}, 125(1), 1--15.
\href{https://doi.org/10.1086/284325}{doi:10.1086/284325}.

\bibitem{garland1993}
Garland, T., Jr., Dickerman, A. W., Janis, C. M., and Jones, J. A. (1993).
Phylogenetic analysis of covariance by computer simulation.
\textit{Systematic Biology}, 42(3), 265--292.
\href{https://doi.org/10.1093/sysbio/42.3.265}{doi:10.1093/sysbio/42.3.265}.

\bibitem{polly2001}
Polly, P. D. (2001).
Paleontology and the comparative method: ancestral node reconstructions versus
observed node values.
\textit{The American Naturalist}, 157(6), 596--609.
\href{https://doi.org/10.1086/320622}{doi:10.1086/320622}.

\bibitem{boettiger2012}
Boettiger, C., Coop, G., and Ralph, P. (2012).
Is your phylogeny informative? Measuring the power of comparative methods.
\textit{Evolution}, 66(7), 2240--2251.
\href{https://doi.org/10.1111/j.1558-5646.2011.01574.x}%
{doi:10.1111/j.1558-5646.2011.01574.x}.

\bibitem{freckleton2012}
Freckleton, R. P. (2012).
Fast likelihood calculations for comparative analyses.
\textit{Methods in Ecology and Evolution}, 3(5), 940--947.
\href{https://doi.org/10.1111/j.2041-210X.2012.00220.x}%
{doi:10.1111/j.2041-210X.2012.00220.x}.

\bibitem{pennell2015}
Pennell, M. W., FitzJohn, R. G., Cornwell, W. K., and Harmon, L. J. (2015).
Model adequacy and the macroevolution of angiosperm functional traits.
\textit{The American Naturalist}, 186(2), E33--E50.
\href{https://doi.org/10.1086/682022}{doi:10.1086/682022}.

\end{thebibliography}

\end{document}
